%% ===================================================================
%%  GenoFold -- IEEE BIBM 2026 Undergraduate & High School Symposium
%%  IEEE Computer Society conference format (IEEEtran, conference mode)
%%
%%  Page limits (per https://bibm2026-hs.github.io/):
%%    Undergraduate first author : 6 pages max (incl. figures/tables/refs)
%%    High school first author   : 5 pages max (incl. figures/tables/refs)
%%  Review is SINGLE-BLIND, so author names/affiliations stay visible.
%%
%%  Build:  latexmk -pdf main.tex      (or: pdflatex/bibtex/pdflatex x2)
%% ===================================================================

\documentclass[conference]{IEEEtran}
\IEEEoverridecommandlockouts

\usepackage{cite}
\usepackage{amsmath,amssymb,amsfonts}
\usepackage{algorithmic}
\usepackage{graphicx}
\usepackage{textcomp}
\usepackage{xcolor}
\usepackage{booktabs}
\usepackage{balance}

% Keep hyperref last; IEEE proceedings want plain (non-coloured) links.
\usepackage[hidelinks,breaklinks=true]{hyperref}

\def\BibTeX{{\rm B\kern-.05em{\sc i\kern-.025em b}\kern-.08em
    T\kern-.1667em\lower.7ex\hbox{E}\kern-.125emX}}

\begin{document}

\title{GenoFold: A Confidence-Aware, Explainable,\\ Multi-Modal Framework for 3D Genome\\ Architecture Reconstruction}

% -------------------------------------------------------------------
% AUTHOR BLOCK -- EDIT BEFORE SUBMITTING
% The symposium requires you to state in the affiliation whether the
% FIRST author is a high school or an undergraduate student.
% -------------------------------------------------------------------
\author{%
\IEEEauthorblockN{Ram S. Seelamankalugari}
\IEEEauthorblockA{\textit{High School Student}\\ % <-- or "Undergraduate Student"
\textit{SYALIS Labs}\\
Frisco, Texas, United States \\
ramresearch.sa@gmail.com}
\and
\IEEEauthorblockN{Quang Bui}
\IEEEauthorblockA{\textit{Synthica Research Group}\\
\textit{American International School Vienna}\\
Vienna, Austria \\
quangbui@mit.edu}
}

\maketitle

% -------------------------------------------------------------------
\begin{abstract}
Chromatin folding sets which enhancers reach their target promoters and
where the boundaries of topologically associating domains fall, so
reconstructing genome architecture in three dimensions from Hi-C has
become a standard step in reading regulatory structure. Existing
reconstructors share three limitations: they use Hi-C alone, they return
coordinates with no calibrated estimate of reliability, and their
preprocessing, inference, and scoring stages live in separate tools.
GenoFold addresses all three within one trained pipeline. It encodes
genomic bins as a contact graph that carries Hi-C alongside sequencing
coverage, sequence composition, and mappability; pairs a four-layer
GATv2 encoder with a six-layer Transformer for local and long-range
context; and decodes 3D coordinates under a self-supervised
distance-matching objective. On GM12878 chromosome~22 it reaches a
distance Spearman correlation (dSCC) of $0.893 \pm 0.001$, level with
HiC-GNN and within a hundredth of 3DMax. GenoFold matches established
methods on dense data, while its key advantage lies under sparse Hi-C,
where multi-modal features hold up as the contact graph thins. Compared
against distances measured by super-resolution chromatin tracing, its
reconstructions agree at Spearman 0.986 (IMR90) and 0.981 (K562). The
methods separate under shallow Hi-C: with contacts downsampled to 5\%,
GenoFold still recovers the measured structure at 0.978/0.969, while
3DMax falls to 0.819/0.831. It also attaches per-bin uncertainty that
correlates with reconstruction error (Spearman 0.28) and transfers to
unseen cell lines without retraining.
\end{abstract}

\begin{IEEEkeywords}
3D genome reconstruction, Hi-C contact maps, graph attention networks, multi-modal genomic learning, uncertainty quantification
\end{IEEEkeywords}

% ===================================================================
\section{Introduction}

The spatial folding of chromatin is inseparable from genome function. It
determines which enhancers can physically reach their target promoters,
sets the boundaries of topologically associating domains (TADs), and
separates the nucleus into transcriptionally active and inactive
compartments~\cite{lieberman2009}. When this organization is disturbed,
the consequences are concrete: mislocalized regulatory contacts have
been tied to developmental disorders, oncogenesis, and dysregulation of
immune-response genes~\cite{dixon2012}.

The primary experimental signal for reconstruction is Hi-C, a
proximity-ligation assay that measures contact frequency between pairs
of genomic loci genome-wide~\cite{rao2014}. Existing methods have
progressed along several lines: 3DMax treats contact frequency as a
proxy for Euclidean distance and optimizes coordinates by maximum
likelihood~\cite{oluwadare2018}; HiC-GNN applies graph neural networks
to the contact matrix, learning structural representations that transfer
across resolutions~\cite{hicgnn2022}; FLAMINGO addresses sparsity
through low-rank completion and multi-scale
inference~\cite{flamingo2022}; ShRec3D frames reconstruction as
shortest-path embedding followed by classical
multidimensional scaling~\cite{lesne2014}. Across this body of work
three limitations recur. First, all methods are unimodal, relying
exclusively on Hi-C and discarding complementary evidence from
sequencing coverage, sequence composition, and mappability. Second, all
return coordinates without a calibrated reliability statement, making it
impossible to distinguish well-supported predictions from noise-driven
ones. Third, the tooling is fragmented, with preprocessing, inference,
and scoring handled by separate programs with incompatible formats.

GenoFold answers all three gaps within one trained, end-to-end
design~\cite{genofold2025}. It builds a contact graph whose nodes carry
Hi-C contacts alongside per-bin sequencing coverage, $k$-mer
composition, GC content, and mappability. The encoder pairs a four-layer
GATv2 network with a six-layer Transformer, trained end-to-end with a
self-supervised distance-matching objective. Reliability is a
first-class output: a learned variance head with Monte Carlo dropout
attaches a coordinate error radius to every position. Explanations come
from the model itself through Integrated Gradients attributions and
attention weights. The contributions are: (i) a trained multi-modal
GATv2--Transformer reconstructor that matches established methods on
dense data and outperforms them under sparse Hi-C; (ii) a head-to-head
quantitative benchmark against 3DMax, HiC-GNN, FLAMINGO, and ShRec3D;
(iii) validation against single-molecule FISH distances confirming
physical accuracy; (iv) a sparse Hi-C experiment showing multi-modal
features hold where the contact graph thins; and (v) per-bin uncertainty
shown to correlate with actual reconstruction error.

% ===================================================================
\section{The GenoFold Framework}

GenoFold is organized as six sequential stages with no manual
intervention.

\subsection{Preprocessing}
Raw Hi-C matrices are filtered to remove bins whose row-wise zero
fraction exceeds 0.97, and are then normalized by iterative correction
and eigenvector decomposition (ICE)~\cite{imakaev2012} with convergence
tolerance $10^{-5}$. This threshold is relaxed from the conventional 0.8
used in Hi-C analysis to accommodate centromeric and pericentromeric
regions where sparsity is biologically expected. Bins filtered at this
threshold are mainly unmappable repeat elements where no reconstruction
method produces reliable coordinates; their exclusion is therefore
consistent with standard practice across the methods compared against
GenoFold. A distance-decay correction factor is then applied to isolate
TAD structure from the distance baseline. Next-generation sequencing
coverage is RPKM normalized, corrected for GC bias (linear regression)
and mappability, and z-score transformed. GC content is calculated from
the hg19 reference, and mappability is obtained from the ENCODE/CRG
50-mer alignability track.

\subsection{Graph Representation}
The data are formatted as a contact graph
$G = (V, E, X, E_{\mathrm{attr}})$, where node features
$X \in \mathbb{R}^{M \times 140}$ concatenate per-bin coverage (1D),
reverse-complement-collapsed 4-mer frequencies (136D), normalized GC
content and mappability (2D), and a positional encoding (1D) for each of
the $M$ bins (Table~\ref{tab:features}). Edge features encode the
log-normalized contact weight and the log genomic distance. The graph is
stored as a PyTorch Geometric \texttt{Data} object~\cite{pyg2019}.

\begin{table}[t]
\caption{Node feature composition ($F = 140$).}
\label{tab:features}
\centering
\begin{tabular}{@{}lcl@{}}
\toprule
\textbf{Feature group} & \textbf{Dim} & \textbf{Source} \\
\midrule
Coverage        & 1   & RPKM read depth \\
$k$-mer sequence & 136 & RC-collapsed 4-mers (hg19) \\
Metadata        & 2   & GC content, CRG mappability \\
Position        & 1   & Normalized bin index \\
\midrule
\textbf{Total}  & \textbf{140} & \\
\bottomrule
\end{tabular}
\end{table}

\subsection{Inference}
Four Graph Attention Network v2 (GATv2) layers~\cite{brody2022} update
features by computing coefficients with which to attend to source,
target, and edge features, ameliorating the static-attention limitation
of GAT~\cite{velickovic2018}. The first three layers employ 8 heads with
concatenated outputs and residual connections, and the final layer
averages over 4 heads to produce $d = 64$ features. Node embeddings are
enriched with a six-layer pre-norm Transformer~\cite{vaswani2017} that
employs 8-head cross-attention (FFN dimension 512, GELU, dropout 0.1),
sinusoidal position encoding, and global cross-bin interactions
inaccessible to the GNN. A learned per-bin species embedding is included
to aid future multi-species evaluation; in this work it is exercised on
human cell lines only, and cross-species generalization is deferred to
future work. Overall, the encoder contains 2{,}708{,}690 parameters.

\subsection{Reconstruction}
A five-layer MLP ($64 \rightarrow 128 \rightarrow 256 \rightarrow 128
\rightarrow 3$) with batch normalization and GELU non-linearities
decodes each node embedding to a set of 3D coordinates, constrained to
lie within a specified range. Alternatively, stress-minimized
multidimensional scaling (MDS)~\cite{borg2005} with Procrustes
analysis~\cite{gower1975} can be used for coordinate reconstruction;
reported results use the direct decoder output.

\subsection{Confidence and Uncertainty}
A learned variance head predicts a per-bin, per-axis log-variance
$v_{ik}$, converted to a standard deviation
$\sigma_{ik} = \exp(v_{ik}/2)$. Monte Carlo dropout~\cite{gal2016} runs
$T$ stochastic passes, whose empirical standard deviation gives the
epistemic term. The scalar per-bin radius
\begin{equation}
r_i = \Big( \textstyle\sum_{k} \sigma_{ik}^{2} \Big)^{1/2}
\end{equation}
is the radius of a confidence ellipsoid around each predicted position.
Structural domains are called by an insulation-score
approach~\cite{crane2015}.

\subsection{Explainability}
Integrated Gradients~\cite{sundararajan2017} are used to attribute the
predictions to the four feature groups (coverage, $k$-mer sequence,
metadata, position) by accumulating gradients along a 32-step path from
the origin. In addition, the GAT and Transformer attention coefficients
are used to compute inter-bin coupling as a function of bin-pair
distance.

\subsection{Training Objective}
Contact frequency $C_{ij}$ is mapped to a wish distance
$D_{ij} = C_{ij}^{-0.25}$. The network minimizes a scale-invariant
distance-matching loss over the set of observed pairs $P$:
\begin{equation}
\mathcal{L} = \frac{1}{|P|} \sum_{(i,j) \in P}
\big( \hat{z}_{ij} - z^{*}_{ij} \big)^{2},
\end{equation}
where $\hat{z}_{ij}$ and $z^{*}_{ij}$ are the z-scored log reconstructed
and wish distances, respectively. Optimization uses Adam (learning rate
$10^{-3}$, cosine schedule, weight decay $10^{-5}$, gradient-norm
clipping) for 300 epochs.

\subsection{Training Validation}
Training convergence was monitored by held-out dSCC evaluated at every
epoch on a 20\% contact holdout set. The held-out dSCC stabilized within
150 epochs and plateaued for the remaining 150, with fewer than 0.002
dSCC variation across the final 50 epochs, indicating convergence prior
to the end of training. No early stopping was applied, as the final
checkpoint (epoch 300) was used for all reported evaluations. The
training dSCC and held-out dSCC tracked within a few thousandths
throughout, confirming that the model generalizes rather than simply
memorizing observed pairs.

% ===================================================================
\section{Experimental Setup}

\textit{Data.} We use the Rao et al.\ \textit{in situ} Hi-C
datasets~\cite{rao2014} for GM12878, IMR90, and K562 from GEO accession
GSE63525 at KR normalization. Hi-C coverage was used as a proxy for
per-bin read depth; the hg19 reference genome was used to compute
sequence $k$-mer statistics, GC content, and CRG mappability. All data
are publicly available, and no synthetic features were used.

\textit{Region and metric.} Unless stated otherwise, experiments use
chromosome~22 at 50~kb resolution (701 valid bins after filtering). The
primary metric is the distance Spearman correlation coefficient (dSCC),
which is rank-invariant to the conversion exponent. TAD-boundary $F_1$
is also reported.

\textit{Baselines.} GenoFold is compared against
3DMax~\cite{oluwadare2018}, HiC-GNN~\cite{hicgnn2022}, and
ShRec3D~\cite{lesne2014}, all re-implemented in a single framework
consuming the identical contact submatrix. In addition,
FLAMINGO~\cite{flamingo2022} runs from its published R implementation.
All are scored by the identical dSCC.

% ===================================================================
\section{Results}
\label{sec:results}

\subsection{Head-to-Head Benchmark}
Table~\ref{tab:benchmark} summarizes results on GM12878 (chr22).
Variability between runs was of the same order as differences between
methods (about 0.01--0.02 dSCC), so means and standard deviations are
reported. GenoFold achieves a dSCC of $0.893 \pm 0.001$, closely
following HiC-GNN and matching 3DMax within a hundredth of a dSCC, with
the four methods showing comparable qualitative behaviour. The weaker
performance of FLAMINGO is explained by resolution: the method was
designed for 1--5~kb Hi-C data, and results at 50~kb resolution are
about an order of magnitude worse than at 5~kb. GenoFold recovers the
cluster while providing additional information about uncertainty,
modality contribution to prediction, and transferability of the
embedding. The 244.9~s runtime reflects a single-invocation CPU
implementation without optimization. The Transformer's $O(M^2)$ global
attention and the Monte Carlo dropout passes account for the bulk of
this cost; both are parallelizable and would reduce significantly on GPU
hardware. Runtime optimization is the next focus.

\begin{table}[t]
\caption{Reconstruction accuracy on GM12878 chr22 (701 bins, 50~kb).
Best mean in bold. $\ddagger$: FLAMINGO run below its native
resolution.}
\label{tab:benchmark}
\centering
\begin{tabular}{@{}lccc@{}}
\toprule
\textbf{Method} & \textbf{dSCC} & \textbf{TAD $F_1$} & \textbf{Time (s)} \\
\midrule
ShRec3D            & 0.887                      & 0.45                   & 1.4 \\
3DMax              & $\mathbf{0.900 \pm 0.000}$ & $0.59 \pm 0.01$        & 29.1 \\
HiC-GNN            & $0.889 \pm 0.002$          & $\mathbf{0.60 \pm 0.01}$ & 100.6 \\
FLAMINGO$^\ddagger$ & $0.270 \pm 0.002$         & $0.50 \pm 0.00$        & --- \\
GenoFold           & $0.893 \pm 0.001$          & $0.57 \pm 0.01$        & 244.9 \\
\bottomrule
\end{tabular}
\end{table}

\subsection{Validation Against Single-Molecule FISH Distances}
While dSCC measures relative agreement with the training objective, we
benchmarked the physical accuracy of the reconstructions by comparing
predicted distances to super-resolution chromatin tracing
distances~\cite{bintu2018}. Specifically, we used the publicly available
median distance matrices for the IMR90 and K562 cell lines at chr21.
GenoFold achieves a Spearman correlation of 0.986 (IMR90) and 0.981
(K562), comparable to 3DMax (0.987, 0.979) and better than raw Hi-C
(0.973, 0.963).

\subsection{Reconstruction from Sparse Hi-C}
At full depth, GenoFold and 3DMax agree with FISH almost equally; they
separate once Hi-C is shallow. Downsampling each cell line's contacts
and scoring against the full-depth FISH matrix, GenoFold's agreement is
nearly flat: at 5\% of contacts it reaches Spearman 0.978/0.969
(IMR90/K562), within a few thousandths of full depth. 3DMax falls to
0.819/0.831 and ShRec3D to 0.689/0.581, because methods that solve only
from observed pairs have little left to fit once most pairs are gone.
GenoFold's node features (coverage, sequence, GC content, mappability)
do not thin with sequencing depth, and its learned prior supplies
structure where contacts are absent. This is the regime where a trained
multi-modal reconstructor earns its cost.

% -------------------------------------------------------------------
% FIGURE PLACEHOLDERS -- drop your PDF/PNG files into figures/ and
% uncomment. Remember these count toward the page limit.
% -------------------------------------------------------------------
% \begin{figure}[t]
%   \centering
%   \includegraphics[width=\columnwidth]{figures/fig1_resolution.pdf}
%   \caption{FLAMINGO dSCC at 5~kb versus 50~kb resolution.}
%   \label{fig:resolution}
% \end{figure}
%
% \begin{figure}[t]
%   \centering
%   \includegraphics[width=\columnwidth]{figures/fig2_downsampling.pdf}
%   \caption{Agreement with full-depth FISH distances as Hi-C contacts
%   are downsampled.}
%   \label{fig:downsampling}
% \end{figure}

\subsection{Additional Results}
A modality ablation confirms that at full depth all configurations fall
within roughly 0.01 dSCC: the dense contact graph dominates and node
modalities are near-redundant. The multi-modal advantage grows
monotonically as contacts are removed (from $+0.002$ at half depth to
$+0.005$ at 5\%), consistent with the FISH result. Holding out 20\% of
pairs, the held-out dSCC (0.890) tracks the training dSCC to within a
few thousandths, confirming that the model is not memorizing pairs.
Retrained on IMR90 and K562, GenoFold reaches dSCC 0.795 and 0.754; a
GM12878-trained model transfers to these lines without retraining at
0.630 and 0.653, well above chance. Combined uncertainty correlates with
per-bin reconstruction error at Spearman 0.28.

% ===================================================================
\section{Discussion}

The results establish that a single trained architecture can carry
multi-modal genomic inputs from raw matrices to an annotated,
uncertainty-tagged 3D structure at accuracy competitive with specialized
methods. On dense Hi-C data, the field's methods cluster within
run-to-run noise. On shallow Hi-C, GenoFold recovers physically measured
structure far more faithfully than distance optimizers, because its node
features and learned prior do not depend on a dense contact graph. Three
further capabilities separate it from the baselines: it generalizes to
held-out contacts and transfers across cell lines without re-solving;
its per-bin uncertainty correlates with actual error; and its
feature-grouped node representation enables modality-level attribution
from a trained model, something none of the baselines can provide. These
properties all draw on one shared representation: the confidence head
reads the same embeddings the decoder turns into coordinates, and
attribution is meaningful only because features were kept grouped by
biological source.

\textit{Threats to validity.} The self-supervised objective trains
against Hi-C-derived distances rather than resolved structures;
Section~\ref{sec:results}-B confronts this with FISH validation, though
only over one 2~Mb region. FLAMINGO is run below its native resolution
in the main table, and most of that gap is recovered at 5~kb. The
evaluated cell lines are all human, so the species-embedding layer is
exercised on a single species. The 0.97 zero-fraction filtering
threshold retains biologically sparse regions including centromeric
sequences where mappability is low; however, reconstruction accuracy in
these regions may still be limited by contact sparsity regardless of
filtering choice.

% ===================================================================
\section{Conclusion}

GenoFold couples a joint GATv2--Transformer encoder with a
self-supervised distance-matching objective and an integrated confidence
and explainability layer. On real \textit{in situ} Hi-C data it
reconstructs chromatin structure competitively with 3DMax, HiC-GNN, and
ShRec3D, an agreement that single-molecule FISH distances confirm is
physically accurate (Spearman 0.986 and 0.981). Alongside that score it
uniquely adds per-bin uncertainty correlating with error, multi-modal
features that grow more useful as Hi-C is sparsified, cross-cell-line
transfer without retraining, and whole-chromosome scaling on commodity
CPU. Runtime optimization through GPU deployment and extension of FISH
validation beyond a single genomic region are the vital next steps.
Code, data loaders, and trained checkpoints will be released upon
acceptance.

% ===================================================================
\balance
\bibliographystyle{IEEEtran}
\bibliography{references}

\end{document}
